%%%%%%%%%%%%%%%%%%%%%%%%%%%%%%%%%%%%%%%%%%%%%%%%%%%%%%%%%%%%%%%%%%
\section{Conclusions}
\label{sec:conclusion}
%%%%%%%%%%%%%%%%%%%%%%%%%%%%%%%%%%%%%%%%%%%%%%%%%%%%%%%%%%%%%%%%%%

\noindent
This paper presents a two-way actuator-line coupling between a geometrically exact Simo--Reissner beam, discretised by a cell-centred finite volume method on a one-dimensional centreline mesh, and an incompressible finite volume flow solver. The fluid velocity is sampled at a prescribed offset upstream of each beam control volume, a quasi-steady Morison-type drag closure evaluates the hydrodynamic line force, and the equal and opposite reaction is projected back into the fluid momentum equation through a normalised Gaussian kernel. Because the two subsystems share one discretisation, one mesh and field infrastructure, one time loop, and one domain decomposition, the coupling reduces to two transfer operators between non-matching meshes rather than an interface between two independently developed codes. The framework was implemented in OpenFOAM, verified against two channel-flow benchmarks, and validated against a wave-flume experiment on a flexible vegetation stem. The following observations are made from the numerical analyses:

\begin{itemize}

\item The projection operator conserves the transferred resultant force by construction, because normalising the accumulated kernel weight per actuator node makes the discrete volume integral of the applied source equal the total actuator-line load to round-off, whatever the mesh, the kernel width, or the kernel truncation. The same guarantee does not extend to angular momentum, and the associated moment discrepancy scales with the kernel width. Both operators were also made parallel consistent on a domain-decomposed fluid mesh, with a bounding-box rejection test, a seeded cell-walking search, and a global-search fallback reducing the per-step search cost to $\mathcal{O}(N_{b})$ in the steady state, and two reduction-broadcast pairs as the only collective communication required.

\item In the steady confined channel benchmark, with the normal drag coefficient fixed at $C_{dn} = 26.73$ by a separate steady CFD calculation on the equivalent rigid plate, the predicted steady tip deflection was $5.73~\mathrm{mm}$, inside the $4.7$ to $6.0~\mathrm{mm}$ band spanned by three published solutions that resolve the same cantilever as a two-dimensional continuum. Six sensitivity studies established that the result is converged in the discretisation sense, the beam mesh, the time step, and the three resolved fluid meshes each changing it by less than $0.5\%$, with the structural Newton tolerance the only remaining sensitivity.

\item The kernel width and the upstream sampling distance were identified as the two operator settings with a physical rather than a purely numerical influence on the response, changing the steady tip deflection by $20\%$ and $31\%$ respectively across their admissible ranges. In a confined channel the two act together as an effective calibration of the quasi-steady closure to the local blockage. The pair selected in the channel benchmark was transferred unchanged to the leaflet benchmark, and the same relative kernel-sizing rule was applied in the wave flume, so what the later cases test is the transferred settings rather than the settings themselves.

\item For the leaflet in a sinusoidal channel flow, the drag coefficient was fixed by two steady CFD calibrations, and the conditions under which the in-channel coefficient can be identified directly with the closure input were derived analytically and confirmed to within $0.6\%$ by a probe-line measurement. Across the four Reynolds and Strouhal combinations of the reference study, the forcing period was reproduced exactly at both Strouhal numbers, the peak tip excursion was over-predicted at a Strouhal number of $1$ at both Reynolds numbers, and it was under-predicted by $30\%$ at a Strouhal number of $0.5$. That under-prediction was traced to the absence of the added-mass and history terms in the closure rather than to the calibration, since an unsteady recalculation of the rigid-plate drag at the two Keulegan--Carpenter numbers returned $C_{d} \approx 107$ in both.

\item In the wave-flume validation the drag coefficient was taken from the fits published alongside the experiment and the operator settings were inherited from the channel benchmarks. The simulated stem motion reproduced the measured mode shape, the single-mode character of the response, the intra-wave asymmetry, the distributed force along the stem, and the base force, and the positive tip excursion matched the measurement to within $1\%$ at the upper published coefficient. The oscillation range was recovered to $65$ to $89\%$ at the primary condition across the propagated coefficient range, and to $71$ to $83\%$ across three wave conditions spanning a threefold change in Keulegan--Carpenter number and tip excursions up to a third of the stem length.

\item A quasi-static reconstruction test, in which the instantaneous actuator-line load was applied statically to an analytic cantilever of the same bending stiffness, reproduced the coupled response to $97$ to $98\%$ in amplitude and exactly in phase at every wave condition. The structural model, its discretisation, and the transmission of the load through the partitioned coupling therefore add no dynamics and no lag of their own, which places the residual $15$ to $20\%$ amplitude deficit on the load side of the chain. That deficit is attributed to the omitted fluid inertia rather than to an incorrect coefficient, since closing it would require a coefficient about a quarter above the largest value the experiment supports.

\end{itemize}

\noindent
Several limitations bound these results. The quasi-steady drag-only closure is the leading one, and it requires a coefficient that must be re-established whenever the Reynolds number or the blockage changes, which no preparatory calculation can supply in a flow that never reaches a steady state. The kernel width and the sampling distance were selected in the channel benchmark with the reference result in view, so the agreement reported there is not a fully blind prediction. The reference data themselves limit the precision of the validation, since the published wave kinematics describe waves exceeding their nominal targets by as much as $17\%$, one condition admits a weak reflected wave, and the blade carries a documented manufacturing pre-tension. The validation also covers a single isolated stem, and a drag-only line force cannot generate vortex-induced lift, so vortex-induced vibration of cables and risers lies outside the reach of this closure regardless of the coupling scheme.

\noindent
The single change that would remove the dominant error is the addition of the Morison inertia term of Eq.~\eqref{eq:fext_addedmass} to the line force, treating the beam-acceleration part on the diagonal of the structural system, which is also what stabilises partitioned coupling at density ratios near unity \citep{causin_addedmass_2005,forster_artificial_2007}. The two deficits quantified above are the targets for that extension. Replacing the single-cell sample with a kernel-weighted volume average \citep{churchfield_advanced_2017,muscari_effective_2024} or a Lagrangian-averaged sample \citep{xie_actuatorline_2021} would then make the sampled velocity smooth in the beam position, which would remove the Newton-loop oscillation that currently forces the interface load to be frozen within the time step and would eliminate the sampling distance as a free parameter. The cost of doing so is a self-induced-deficit correction of the kind developed for lift lines \citep{martinez_tossas_filtered_2019,meyerforsting_vortexbased_2019,kleine_noniterative_2023}. With smooth sampling in place, an implicit split of the drag source under Aitken dynamic relaxation \citep{kuttler_fixedpoint_2008} or quasi-Newton acceleration \citep{degroote_performance_2010} would convert the present staggered scheme into a strongly coupled iteration, and an alternative route makes the added mass implicit in the rod update instead \citep{tschisgale_immersed_2020,tschisgale_large_2021}. Restricting the projection sweep to a band of cells around the beam polyline is the remaining prerequisite for the canopy and net configurations the method is intended for, and the per-stem accuracy quantified here is what a canopy result would inherit.
